\corrigesChapitre{13}{Algorithmique, logique et raisonnement}
\begin{corrige}{1}{Affectation}
\begin{questions}
\item $u=7$.
\item Le signe d’affectation remplace une ancienne valeur par une nouvelle ; il n’exprime pas une égalité simultanée.
\item $2\times7+1=15$.
\end{questions}
\end{corrige}
\begin{corrige}{2}{Liste}
\begin{questions}
\item $[0,1,3,7]$.
\item $4$, puisque la valeur initiale est conservée avant les trois mises à jour.
\item $3$, qui est $u_2$.
\end{questions}
\end{corrige}
\begin{corrige}{3}{Ensembles}
\begin{questions}
\item $A\cap B=\{3\}$ ; $A\cup B=\{1,2,3,4\}$.
\item $\{4,5\}$.
\item Non : $1\in A$ mais $1\notin B$.
\end{questions}
\end{corrige}
\begin{corrige}{4}{Nier une proposition}
\begin{questions}
\item « $x\leq2$ ».
\item « $x\geq0$ et $x\leq4$ ».
\item « $x<1$ ou $x>3$ ». La négation échange « et » et « ou » et nie chaque condition.
\end{questions}
\end{corrige}
\begin{corrige}{5}{Finir une boucle}
\begin{questions}
\item $7$ : $2^6=64<100\leq128=2^7$.
\item $128$.
\item La boucle ferait encore une mise à jour quand $u=128$ ; elle s’arrêterait en $256$, après $8$ mises à jour.
\end{questions}
\end{corrige}
\begin{corrige}{6}{Une fonction de moyenne}
\begin{questions}
\item \code{def moyenne(valeurs): return sum(valeurs)/len(valeurs)}.
\item La liste doit être non vide et ses éléments numériques ; sinon la division ou la somme n’a pas le sens voulu.
\item $(2+4+9)/3=5$.
\end{questions}
\end{corrige}
\begin{corrige}{7}{Implication}
\begin{questions}
\item Oui : $x>2$ est positif et le carré est strictement croissant sur $[0,+\infty[$.
\item « $x^2>4$ entraîne $x>2$ » est fausse sur $\R$, avec $x=-3$.
\item Sur $[0,+\infty[$, $x>2$ équivaut à $x^2>4$.
\end{questions}
\end{corrige}
\begin{corrige}{8}{Tester un programme de cumul}
\begin{questions}
\item $1$ mathématiquement, car toutes les valeurs possibles sont incluses.
\item $(1-p)^n$.
\item Les arrondis successifs sur des nombres flottants peuvent créer une petite erreur ; on contrôle avec une tolérance adaptée.
\end{questions}
\end{corrige}
\begin{corrige}{9}{Contre-exemple}
\begin{questions}
\item $u_n=1/(n+1)$ est positive et strictement décroissante.
\item $f(x)=x^3$ a $f\,\!^{\prime}(0)=0$ sans extremum en zéro.
\item Une affirmation universelle prétend couvrir tous les cas ; un cas contraire contredit cette prétention.
\end{questions}
\end{corrige}
\begin{corrige}{10}{Raisonner sur un intervalle}
\begin{questions}
\item $x-1>0$ et $3-x>0$, donc $1<x<3$.
\item La stricte croissance de $\ln$ donne $x-1=3-x$, donc $x=2$, qui appartient au domaine.
\item Annoncer le domaine, justifier l’injectivité du logarithme, résoudre puis vérifier l’appartenance. On conclut par l’ensemble $\{2\}$.
\end{questions}
\end{corrige}
\begin{corrige}{11}{Somme : une erreur de borne}
\begin{questions}
\item $0+1+\cdots+(n-1)$ pour $n\geq1$.
\item Utiliser \code{sum(range(1, n+1))}.
\item Le premier rend $10$ ; le second $15$. La borne finale de \code{range} est exclue.
\end{questions}
\end{corrige}
\begin{corrige}{12}{Un oral argumenté}
\begin{questions}
\item Par exemple le modèle fictif du chapitre 8 : proportion initiale $2\,\%$, sensibilité $95\,\%$, spécificité $95\,\%$.
\item Présenter un arbre ou les effectifs sur $10\,000$ personnes : $190$ vrais positifs et $490$ faux positifs ; Bayes donne $19/68\simeq27{,}94\,\%$.
\item La sensibilité et la valeur prédictive positive sont deux conditionnements différents. La conclusion dépend des données fictives et des hypothèses du modèle, pas d’un conseil individuel.
\end{questions}
\end{corrige}
